% from ega0-19-fr.tex
\oldpage[0\textsubscript{IV-1}]{75}

  \item[b)] L'homomorphisme
  $u_0:M/\mathfrak QM\to N/\mathfrak QN$ déduit de $u$ par passage aux
  quotients est inversible à gauche.
\end{enumerate}
\end{proposition}

\begin{proof}
On a vu (19.1.7) que la condition a) équivaut à dire que $u_\lambda$ est
inversible à gauche pour tout $\lambda$~; comme $\mathfrak Q$ est un idéal
ouvert, donc contient un $\mathfrak J_\lambda$, on en déduit aussitôt que
$u_0$ est inversible à gauche. Pour montrer inversement que b) entraîne a),
notons que pour tout $\lambda$, $\mathfrak Q/\mathfrak J_\lambda$ est par
hypothèse un idéal nilpotent de $A/\mathfrak J_\lambda$. Notre assertion
résultera de la proposition suivante.
\end{proof}

\begin{proposition}[19.1.10]
\label{0.19.1.10-fr}
Soient $A$ un anneau, $M$, $N$ deux $A$-modules, $N$ étant projectif,
$u:M\to N$ un $A$-homomorphisme. Soit $\mathfrak J$ un idéal de $A$~; on
suppose vérifiée l'une des conditions suivantes~:
\begin{enumerate}
  \item[(i)] $\mathfrak J$ est nilpotent.
  \item[(ii)] $\mathfrak J$ est contenu dans le radical de $A$ et $M$ est de
  type fini.
\end{enumerate}
Alors, pour que $u$ soit inversible à gauche, il faut et il suffit que
l'homomorphisme
$u_0:M/\mathfrak JM\to N/\mathfrak JN$ de $(A/\mathfrak J)$-modules, déduit
de $u$ par passage aux quotients, soit inversible à gauche.
\end{proposition}

\begin{proof}
La condition étant évidemment nécessaire, prouvons qu'elle est suffisante.
Soit $v_0$ un inverse à gauche de $u_0$~; l'homomorphisme composé
$N\to N/\mathfrak JN\xrightarrow{v_0}M/\mathfrak JM$ se factorise en
$N\xrightarrow{v}M\to M/\mathfrak JM$ puisque $N$ est projectif~; alors
$w=v\circ u$ est un endomorphisme de $M$ tel que l'endomorphisme $w_0$ de
$M/\mathfrak JM$ déduit par passage aux quotients soit l'identité, et il suffit
de prouver que $w$ est lui-même bijectif (car alors $w^{-1}\circ v$ sera un
inverse à gauche de $u$). Distinguons maintenant les deux cas.

(i) Pour tout $n$ on a le diagramme commutatif
\[
  \xymatrix{
    (\mathfrak J^n/\mathfrak J^{n+1})\otimes_{A/\mathfrak J}
      (M/\mathfrak JM)\ar[r]\ar[d]_{1\otimes\operatorname{gr}^0(w)} &
      \mathfrak J^nM/\mathfrak J^{n+1}M\ar[d]^{\operatorname{gr}^n(w)}\\
    (\mathfrak J^n/\mathfrak J^{n+1})\otimes_{A/\mathfrak J}
      (M/\mathfrak JM)\ar[r] & \mathfrak J^nM/\mathfrak J^{n+1}M
  }
\]
où les flèches horizontales sont \emph{surjectives}, et comme
$\operatorname{gr}^0(w)=w_0$ est l'identité, il en est de même de
$\operatorname{gr}^n(w)$ qui, \emph{a fortiori}, est bijectif. La filtration
$\mathfrak J$-préadique sur $M$ étant finie puisque $\mathfrak J$ est
nilpotent, on en conclut que $w$ est bijectif (Bourbaki, \emph{Alg. comm.},
chap.~III, §~2, n\textsuperscript{o}~8, cor.~3 du th.~1).

(ii) Il suffit de montrer que pour tout idéal maximal $\mathfrak m$ de $A$,
l'endomorphisme $w_{\mathfrak m}$ de $M_{\mathfrak m}$ est bijectif
(Bourbaki, \emph{Alg. comm.}, chap.~II, §~3, n\textsuperscript{o}~3, th.~1)
et comme $\mathfrak JA_{\mathfrak m}\subset\mathfrak mA_{\mathfrak m}$, et
$A_{\mathfrak m}/\mathfrak JA_{\mathfrak m}=(A/\mathfrak J)_{\mathfrak m}$,
on est ramené à prouver la proposition lorsque $A$ est un anneau local. En
outre, on peut supposer que $\mathfrak J$ est l'idéal maximal
$\mathfrak J_0$ de $A$, car si $u_0$ est inversible

